\subsection{Coincidencia de Patrones}

\subsubsection{Regla: patrón de tipo en lugar de \texttt{as} con comprobación de nulo}

\textbf{Regla:} Para comprobar el tipo de un valor y usarlo convertido se emplea el patrón \texttt{is Tipo nombre}, o \texttt{is not Tipo nombre} cuando forma parte de una cláusula de guarda (Sección 4.2). Prohibido usar \texttt{as} seguido de una comprobación de nulo o una conversión explícita después de \texttt{is}. La variable que introduce el patrón sigue las reglas de nombrado de variables (Sección 2.1) (Microsoft, 2026i).

\textbf{Descripción:} El patrón comprueba el tipo y declara la variable en una sola expresión, y el compilador solo la considera asignada en la rama donde el patrón se cumple, por lo que no puede usarse por error cuando la comprobación falla.

\textbf{Código correcto:}
\begin{lstlisting}[style=csharpstyle]
public void ApplyEffect(Tile tile)
{
    if (tile is not TrapTile trapTile)
    {
        return;
    }

    trapTile.Trigger();
}
\end{lstlisting}

\textbf{Código incorrecto:}
\begin{lstlisting}[style=csharpstyle]
public void ApplyEffect(Tile tile)
{
    TrapTile trapTile = tile as TrapTile;

    if (trapTile == null)
    {
        return;
    }

    trapTile.Trigger();
}
\end{lstlisting}
